% GENERATED FILE. Edit canonical lessons, then run npm run generate:books.
% canonical-source-hash: fnv1a64:f70d86e42b322f46
% canonical-lessons: RU-C256-professor, RU-C256-zhurnalist, RU-C256-khudozhnik, RU-C256-muzykant, RU-C256-aktyor

\chapter{Professor, Journalist, Artist, Musician, Actor}
\label{ch:ru-professor-journalist-artist-musician-actor}
\hlchaptermodality{256}

\begin{chapteropening}
\textbf{By the end of this chapter:} \emph{I can say a professor, a journalist, an artist, a musician and an actor.}

Complete the last lesson of chapter 256: I can say a professor, a journalist, an artist, a musician and an actor.
\end{chapteropening}

\section[\texorpdfstring{\ru{профессор} (proféssor)}{профессор (proféssor)}]{\ru{профессор} (proféssor) — a professor}

\label{lesson:RU-C256-professor}



\begin{quote}
Before the new one: say the Russian for a message, then the Russian for a conversation.
\end{quote}

\subsection*{You'll want to know: \ru{профессор}}
\textbf{\ru{профессор}} (\emph{proféssor}) — \textquotedblleft{}a professor\textquotedblright{}.

\begin{cousinweb}
\textbf{\ru{профессор}} — from 1720; it came through German \emph{Professor} from Latin \emph{professor}, \textquotedblleft{}a public teacher\textquotedblright{}, from \emph{profitērī}, \textquotedblleft{}to declare openly\textquotedblright{} (Vasmer).
\end{cousinweb}

\begin{grammarlens}[title={What you've built}]
One more word for this chapter.
\end{grammarlens}

\subsection*{Your turn}
\begin{itemize}
\raggedright
  \item \emph{Say it:} \emph{proféssor}
  \item \emph{Say it:} \emph{proféssor}, once more
  \item \emph{Say it:} say \emph{proféssor}
  \item \emph{Recall:} say the Russian for a message, then the Russian for a conversation, then say \emph{proféssor} again
\end{itemize}

\begin{tcolorbox}[breakable,colback=teal!4,colframe=teal!35!black,title={Before you move on}]
What does \ru{профессор} mean? (\textquotedblleft{}A professor\textquotedblright{}.) Say it once more.
\end{tcolorbox}
\section[\texorpdfstring{\ru{журналист} (zhurnalíst)}{журналист (zhurnalíst)}]{\ru{журналист} (zhurnalíst) — a journalist}

\label{lesson:RU-C256-zhurnalist}



\begin{quote}
Before the new one: say the Russian for a conversation, then the Russian for a professor.
\end{quote}

\subsection*{You'll want to know: \ru{журналист}}
\textbf{\ru{журналист}} (\emph{zhurnalíst}) — \textquotedblleft{}a journalist\textquotedblright{}.

\begin{cousinweb}
\textbf{\ru{журналист}} is \emph{\ru{журнал}}, \textquotedblleft{}a journal\textquotedblright{}, with \emph{-\ru{ист}}. \emph{\ru{журнал}} comes from French \emph{journal}, Medieval Latin \emph{diurnālis}, \textquotedblleft{}daily\textquotedblright{} (Vasmer).
\end{cousinweb}

\begin{grammarlens}[title={What you've built}]
One more word for this chapter.
\end{grammarlens}

\subsection*{Your turn}
\begin{itemize}
\raggedright
  \item \emph{Say it:} \emph{zhurnalíst}
  \item \emph{Say it:} \emph{zhurnalíst}, once more
  \item \emph{Say it:} say \emph{zhurnalíst}
  \item \emph{Recall:} say the Russian for a conversation, then the Russian for a professor, then say \emph{zhurnalíst} again
\end{itemize}

\begin{tcolorbox}[breakable,colback=teal!4,colframe=teal!35!black,title={Before you move on}]
What does \ru{журналист} mean? (\textquotedblleft{}A journalist\textquotedblright{}.) Say it once more.
\end{tcolorbox}
\section[\texorpdfstring{\ru{художник} (khudózhnik)}{художник (khudózhnik)}]{\ru{художник} (khudózhnik) — an artist, a painter}

\label{lesson:RU-C256-khudozhnik}



\begin{quote}
Before the new one: say the Russian for a professor, then the Russian for a journalist.
\end{quote}

\subsection*{You'll want to know: \ru{художник}}
\textbf{\ru{художник}} (\emph{khudózhnik}) — \textquotedblleft{}an artist, a painter\textquotedblright{}.

\begin{cousinweb}
\textbf{\ru{художник}} comes from Old Russian \emph{xudogŭ}, \textquotedblleft{}skilled\textquotedblright{}, from \emph{*xǫdogŭ}, borrowed from a Gothic \emph{*handags}, \textquotedblleft{}handy\textquotedblright{}, from \emph{handus}, \textquotedblleft{}hand\textquotedblright{} (Vasmer). English \emph{hand} is a cousin of the Gothic word.
\end{cousinweb}

\begin{grammarlens}[title={What you've built}]
One more word for this chapter.
\end{grammarlens}

\subsection*{Your turn}
\begin{itemize}
\raggedright
  \item \emph{Say it:} \emph{khudózhnik}
  \item \emph{Say it:} \emph{khudózhnik}, once more
  \item \emph{Say it:} say \emph{khudózhnik}
  \item \emph{Recall:} say the Russian for a professor, then the Russian for a journalist, then say \emph{khudózhnik} again
\end{itemize}

\begin{tcolorbox}[breakable,colback=teal!4,colframe=teal!35!black,title={Before you move on}]
What does \ru{художник} mean? (\textquotedblleft{}An artist, a painter\textquotedblright{}.) Say it once more.
\end{tcolorbox}
\section[\texorpdfstring{\ru{музыкант} (muzykánt)}{музыкант (muzykánt)}]{\ru{музыкант} (muzykánt) — a musician}

\label{lesson:RU-C256-muzykant}



\begin{quote}
Before the new one: say the Russian for a journalist, then the Russian for an artist.
\end{quote}

\subsection*{You'll want to know: \ru{музыкант}}
\textbf{\ru{музыкант}} (\emph{muzykánt}) — \textquotedblleft{}a musician\textquotedblright{}.

\begin{cousinweb}
\textbf{\ru{музыкант}} — from Peter the Great's time; it comes from German \emph{Musikant}, from Latin \emph{mūsicus}, \textquotedblleft{}a musician\textquotedblright{} (Vasmer).
\end{cousinweb}

\begin{grammarlens}[title={What you've built}]
One more word for this chapter.
\end{grammarlens}

\subsection*{Your turn}
\begin{itemize}
\raggedright
  \item \emph{Say it:} \emph{muzykánt}
  \item \emph{Say it:} \emph{muzykánt}, once more
  \item \emph{Say it:} say \emph{muzykánt}
  \item \emph{Recall:} say the Russian for a journalist, then the Russian for an artist, then say \emph{muzykánt} again
\end{itemize}

\begin{tcolorbox}[breakable,colback=teal!4,colframe=teal!35!black,title={Before you move on}]
What does \ru{музыкант} mean? (\textquotedblleft{}A musician\textquotedblright{}.) Say it once more.
\end{tcolorbox}
\section[\texorpdfstring{\ru{актёр} (aktyór)}{актёр (aktyór)}]{\ru{актёр} (aktyór) — an actor}

\label{lesson:RU-C256-aktyor}



\begin{quote}
Before the new one: say the Russian for an artist, then the Russian for a musician.
\end{quote}

\subsection*{You'll want to know: \ru{актёр}}
\textbf{\ru{актёр}} (\emph{aktyór}) — \textquotedblleft{}an actor\textquotedblright{}.

\begin{cousinweb}
\textbf{\ru{актёр}} comes from French \emph{acteur} (Vasmer), from Latin \emph{āctor}, \textquotedblleft{}a doer, a performer\textquotedblright{}.
\end{cousinweb}

\begin{grammarlens}[title={What you've built}]
One more word for this chapter.
\end{grammarlens}

\subsection*{Your turn}
\begin{itemize}
\raggedright
  \item \emph{Say it:} \emph{aktyór}
  \item \emph{Say it:} \emph{aktyór}, once more
  \item \emph{Say it:} say \emph{aktyór}
  \item \emph{Recall:} say the Russian for an artist, then the Russian for a musician, then say \emph{aktyór} again
\end{itemize}

\begin{tcolorbox}[breakable,colback=teal!4,colframe=teal!35!black,title={Before you move on}]
What does \ru{актёр} mean? (\textquotedblleft{}An actor\textquotedblright{}.) Say it once more.
\end{tcolorbox}
